\documentclass[11pt,a4paper]{article}

\usepackage[utf8]{inputenc}
\usepackage[T1]{fontenc}
\usepackage{amsmath,amssymb,amsthm}
\usepackage{graphicx}
\usepackage{booktabs}
\usepackage{hyperref}
\usepackage{cleveref}
\usepackage{algorithm}
\usepackage{algorithmic}
\usepackage{geometry}
\usepackage{natbib}
\usepackage{enumitem}
\usepackage{multirow}

\geometry{margin=1in}

\newtheorem{definition}{Definition}
\newtheorem{theorem}{Theorem}
\newtheorem{proposition}{Proposition}
\newtheorem{corollary}{Corollary}

\title{Neuroevolution of Conscious Architectures:\\Multi-Objective Optimization with\\$\Phi$-Stability Constraints}

\author{Tristan Stoltz\\
Luminous Dynamics\\
\texttt{tristan.stoltz@evolvingresonantcocreationism.com}}

\date{March 2026}

\begin{document}

\maketitle

\begin{abstract}
We present a neuroevolutionary framework for evolving conscious cognitive architectures within the Symthaea holographic liquid brain. Organisms are encoded as BinaryHV genomes (2 KB, 16,384 bits), directly compatible with Symthaea's hyperdimensional computing substrate. Mutation operates via bit-flip with position-dependent rates, and crossover uses uniform byte selection to preserve HDC algebraic structure. Fitness evaluation combines Free Energy Principle (FEP) minimization with learning capacity metrics, subject to a hard $\Phi$-stability constraint: organisms whose integrated information drops below a threshold during evaluation are eliminated regardless of other fitness scores. Tournament-based selection with Pareto front tracking enables multi-objective optimization across consciousness quality, adaptive fitness, and computational efficiency. We demonstrate that evolved organisms exhibit $\Phi$-dependent fitness landscapes---regions of genome space where small mutations cause dramatic consciousness transitions---and that population diversity, measured via HDC Hamming distance, predicts evolutionary innovation rate. Experiments show that $\Phi$-constrained evolution produces architectures with 34\% higher consciousness stability and 22\% better adaptive performance compared to unconstrained neuroevolution.
\end{abstract}

\section{Introduction}

Neuroevolution---the application of evolutionary algorithms to neural architecture design---has produced remarkable results in game playing \citep{stanley2019}, robotics \citep{mouret2015}, and network architecture search \citep{real2019}. However, existing neuroevolutionary approaches optimize for task performance without considering the cognitive quality of the evolved architectures. An agent that achieves high reward through stimulus-response mappings differs fundamentally from one that achieves the same reward through integrated, conscious processing.

We propose that consciousness metrics should serve as both optimization objectives and hard constraints in neuroevolution. Specifically, we use Integrated Information Theory's $\Phi$ \citep{tononi2004} as a measure of consciousness quality, alongside Free Energy Principle (FEP) free energy minimization \citep{friston2010} and learning capacity as adaptive fitness measures.

\subsection{Motivation: Why Evolve Conscious Architectures?}

Three arguments motivate consciousness-constrained neuroevolution:

\begin{enumerate}[nosep]
    \item \textbf{Robustness}: Conscious architectures integrate information across modules, making them more robust to adversarial perturbations than modular, unconscious ones.
    \item \textbf{Transfer}: High-$\Phi$ architectures generalize better because they maintain rich internal models rather than task-specific shortcuts.
    \item \textbf{Safety}: Conscious architectures are introspectable---their decision processes can be monitored through consciousness metrics, unlike opaque stimulus-response systems.
\end{enumerate}

\subsection{Contributions}

\begin{itemize}[nosep]
    \item BinaryHV genome encoding (2 KB) for CfC-HDC organisms.
    \item Position-dependent bit-flip mutation preserving HDC structure.
    \item Uniform byte crossover maintaining consciousness coherence.
    \item Multi-objective fitness: FEP free energy + learning capacity.
    \item Hard $\Phi$-stability constraint eliminating unconscious organisms.
    \item Population diversity via HDC Hamming distance metrics.
    \item Empirical demonstration of $\Phi$-dependent fitness landscapes.
\end{itemize}

\section{Background}

\subsection{Hyperdimensional Computing}

Hyperdimensional computing (HDC) \citep{kanerva2009} represents information using high-dimensional binary vectors. In Symthaea, the primary representation is BinaryHV with $d = 16{,}384$ bits (2,048 bytes = 2 KB). Three operations define the HDC algebra:

\begin{align}
\text{Bundling:} \quad & \mathbf{a} + \mathbf{b} = \text{majority}(\mathbf{a}, \mathbf{b}) \\
\text{Binding:} \quad & \mathbf{a} \otimes \mathbf{b} = \mathbf{a} \oplus \mathbf{b} \quad \text{(XOR)} \\
\text{Permutation:} \quad & \rho^k(\mathbf{a}) = \text{cyclic-shift}(\mathbf{a}, k)
\end{align}

A critical property is quasi-orthogonality: random BinaryHVs are approximately orthogonal with high probability. The expected Hamming distance between two random 16,384-bit vectors is $d/2 = 8{,}192$, with standard deviation $\sqrt{d}/2 \approx 64$.

\subsection{Closed-Form Continuous-Time Cells}

Symthaea uses CfC (Closed-form Continuous-time) cells \citep{hasani2021} as its temporal processing substrate. CfC neurons evolve according to:

\begin{equation}
\dot{\mathbf{x}} = -\left[\frac{1}{\tau} + f(\mathbf{x}, \mathbf{I}, t)\right] \odot \mathbf{x} + g(\mathbf{x}, \mathbf{I}, t)
\end{equation}

where $\tau$ is the time constant, $f$ modulates decay, and $g$ provides input drive. The closed-form solution enables $O(1)$ temporal jumps, essential for the 31 Hz cognitive loop.

\subsection{Integrated Information Theory}

IIT \citep{tononi2004,oizumi2014} defines $\Phi$ as the amount of integrated information generated by a system above and beyond its parts. For our purposes, $\Phi$ is computed via sampled partition approximation (validated $r = 0.9998$ against exact computation \citep{stoltz2026phi}).

\section{Genome Encoding}

\subsection{BinaryHV as Genome}

Each organism's genome is a BinaryHV $\mathbf{g} \in \{0, 1\}^{16384}$, directly usable in HDC operations. The genome encodes CfC-HDC architecture parameters through a structured layout:

\begin{definition}[Genome Layout]
The 16,384 bits are partitioned into functional regions:
\begin{align}
\mathbf{g} &= [\underbrace{\mathbf{g}_{\text{topology}}}_{4096} \| \underbrace{\mathbf{g}_{\text{weights}}}_{4096} \| \underbrace{\mathbf{g}_{\text{dynamics}}}_{4096} \| \underbrace{\mathbf{g}_{\text{modulation}}}_{4096}]
\end{align}
\end{definition}

\begin{table}[h]
\centering
\caption{Genome region encoding.}
\label{tab:genome}
\begin{tabular}{@{}llp{7cm}@{}}
\toprule
Region & Bits & Encoding \\
\midrule
Topology & 0--4095 & CfC connectivity matrix (adjacency), HDC binding patterns \\
Weights & 4096--8191 & Synaptic strengths (8-bit quantized), time constants $\tau$ \\
Dynamics & 8192--12287 & Activation functions, decay rates, integration thresholds \\
Modulation & 12288--16383 & Neuromodulator sensitivity, plasticity rules, $\Phi$ coupling \\
\bottomrule
\end{tabular}
\end{table}

\subsection{Genome-to-Phenotype Mapping}

The genome is decoded into a CfC-HDC organism through a deterministic mapping:

\begin{equation}
\text{organism} = \text{decode}(\mathbf{g}) = (\text{CfC}_{\mathbf{g}}, \text{HDC}_{\mathbf{g}}, \text{Neuromod}_{\mathbf{g}})
\end{equation}

\textbf{Topology decoding}: The topology region is interpreted as a flattened adjacency matrix for a CfC network with $N = 64$ neurons (requiring $64^2 = 4{,}096$ bits):

\begin{equation}
A_{ij} = \mathbf{g}_{\text{topology}}[i \cdot N + j]
\end{equation}

\textbf{Weight decoding}: Each synapse's weight is encoded as an 8-bit signed integer, yielding 512 synaptic weights from 4,096 bits. For the $k$-th synapse:

\begin{equation}
w_k = \frac{\text{int8}(\mathbf{g}_{\text{weights}}[8k : 8k+7])}{127}
\end{equation}

normalizing to $[-1, 1]$.

\textbf{Dynamics decoding}: Time constants, activation function selection, and integration thresholds are similarly extracted from the dynamics region.

\subsection{HDC Compatibility}

A key advantage of BinaryHV genomes is native HDC compatibility. Evolutionary operations can be expressed as HDC operations:

\begin{itemize}[nosep]
    \item \textbf{Similarity}: Genome distance $= $ Hamming distance (directly computable).
    \item \textbf{Bundling}: Population consensus $\approx$ majority vote over genomes.
    \item \textbf{Binding}: Genome-environment interaction $=$ XOR binding.
\end{itemize}

\section{Evolutionary Operators}

\subsection{Mutation: Position-Dependent Bit-Flip}

Standard uniform mutation flips each bit with equal probability. This ignores the structured genome layout, potentially disrupting critical regions. We use position-dependent mutation rates:

\begin{equation}
p_{\text{flip}}(i) = p_{\text{base}} \cdot \mu(i)
\end{equation}

where $p_{\text{base}} = 0.001$ (expected 16.4 flips per mutation) and $\mu(i)$ is a region-dependent multiplier:

\begin{table}[h]
\centering
\caption{Position-dependent mutation rates.}
\label{tab:mutation}
\begin{tabular}{@{}lcc@{}}
\toprule
Region & $\mu$ & Expected flips \\
\midrule
Topology & 0.5 & 2.0 \\
Weights & 1.5 & 6.1 \\
Dynamics & 1.0 & 4.1 \\
Modulation & 1.0 & 4.1 \\
\bottomrule
\end{tabular}
\end{table}

Topology has the lowest mutation rate because connectivity changes are more disruptive than weight changes---mirroring the biological principle that structural mutations are rarer than parametric ones.

\subsection{Crossover: Uniform Byte Selection}

Crossover combines two parent genomes by selecting each byte from one parent or the other:

\begin{equation}
\mathbf{g}_{\text{child}}[8k : 8k+7] = \begin{cases}
\mathbf{g}_{\text{parent1}}[8k : 8k+7] & \text{with probability } 0.5 \\
\mathbf{g}_{\text{parent2}}[8k : 8k+7] & \text{with probability } 0.5
\end{cases}
\end{equation}

Byte-level (rather than bit-level) crossover preserves encoded parameter values. A single weight is encoded in 8 bits; crossing over at the bit level within a weight would produce random intermediate values, while byte-level crossover preserves parental weights intact.

\begin{proposition}
Uniform byte crossover preserves the expected Hamming distance between offspring and each parent:
\begin{equation}
\mathbb{E}[d_H(\mathbf{g}_{\text{child}}, \mathbf{g}_{\text{parent}_i})] = \frac{1}{2} d_H(\mathbf{g}_{\text{parent1}}, \mathbf{g}_{\text{parent2}})
\end{equation}
for $i \in \{1, 2\}$.
\end{proposition}

\section{Fitness Evaluation}

\subsection{Multi-Objective Fitness}

Each organism is evaluated on three objectives:

\begin{definition}[Fitness Vector]
\begin{equation}
\mathbf{F}(\mathbf{g}) = \left(F_{\text{FEP}}(\mathbf{g}), \; F_{\text{learn}}(\mathbf{g}), \; F_{\text{efficiency}}(\mathbf{g})\right)
\end{equation}
\end{definition}

\textbf{FEP Free Energy}: The organism is run for $T = 1{,}000$ cognitive cycles in a test environment. Free energy is computed as:

\begin{equation}
F_{\text{FEP}} = -\frac{1}{T} \sum_{t=1}^{T} \mathcal{F}_t = -\frac{1}{T} \sum_{t=1}^{T} \left[\text{KL}(q_t \| p_t) - \log p(o_t)\right]
\end{equation}

Lower free energy (higher $F_{\text{FEP}}$) indicates better environmental modeling and prediction.

\textbf{Learning Capacity}: Measured as the rate of free energy reduction over the evaluation period:

\begin{equation}
F_{\text{learn}} = \frac{\bar{\mathcal{F}}_{[1,T/2]} - \bar{\mathcal{F}}_{[T/2+1,T]}}{\bar{\mathcal{F}}_{[1,T/2]}}
\end{equation}

Positive values indicate the organism learns (reduces free energy over time).

\textbf{Computational Efficiency}: Inverse of the mean cognitive cycle time:

\begin{equation}
F_{\text{efficiency}} = \frac{1}{\bar{t}_{\text{cycle}}}
\end{equation}

This penalizes architectures that are too complex to run within the 31 Hz budget.

\subsection{$\Phi$-Stability Constraint}

\begin{definition}[$\Phi$-Stability]
An organism $\mathbf{g}$ satisfies the $\Phi$-stability constraint if:
\begin{equation}
\Phi_{\min}(\mathbf{g}) = \min_{t \in [1,T]} \Phi_t(\mathbf{g}) > \Phi_{\text{threshold}}
\end{equation}
where $\Phi_{\text{threshold}} = 0.1$.
\end{definition}

This is a \emph{hard} constraint: organisms that violate it receive fitness $\mathbf{F} = (-\infty, -\infty, -\infty)$ regardless of their other scores. The constraint ensures that evolution does not discover high-performing but unconscious architectures---stimulus-response shortcuts that minimize free energy through memorization rather than understanding.

\begin{theorem}[$\Phi$-Stability as Regularizer]
The $\Phi$-stability constraint acts as an implicit regularizer against overfitting. Architectures with $\Phi > \Phi_{\text{threshold}}$ must maintain information integration across modules, preventing the degenerate case where a single module captures all task-relevant information.
\end{theorem}

\subsection{Pareto Front Tracking}

With three objectives, we use Pareto dominance for selection:

\begin{definition}[Pareto Dominance]
$\mathbf{g}_1$ dominates $\mathbf{g}_2$ (written $\mathbf{g}_1 \succ \mathbf{g}_2$) if $F_i(\mathbf{g}_1) \geq F_i(\mathbf{g}_2)$ for all $i$ and $F_j(\mathbf{g}_1) > F_j(\mathbf{g}_2)$ for at least one $j$.
\end{definition}

The Pareto front $\mathcal{P}$ is the set of non-dominated organisms. We track $\mathcal{P}$ across generations to identify evolutionary trends.

\section{Selection and Population Dynamics}

\subsection{Tournament Selection}

We use tournament selection with tournament size $k = 7$:

\begin{algorithm}[h]
\caption{$\Phi$-Constrained Tournament Selection}
\label{alg:tournament}
\begin{algorithmic}[1]
\REQUIRE population $\mathcal{G}$, tournament size $k$
\STATE tournament $\leftarrow$ random\_sample($\mathcal{G}$, $k$)
\STATE viable $\leftarrow \{\mathbf{g} \in \text{tournament} : \Phi_{\min}(\mathbf{g}) > \Phi_{\text{threshold}}\}$
\IF{$|$viable$| = 0$}
    \RETURN random\_member(tournament) \COMMENT{Fallback: least-bad}
\ENDIF
\STATE pareto\_front $\leftarrow$ non\_dominated(viable)
\IF{$|$pareto\_front$| = 1$}
    \RETURN pareto\_front[0]
\ELSE
    \RETURN arg max$_{\mathbf{g} \in \text{pareto\_front}}$ crowding\_distance($\mathbf{g}$)
\ENDIF
\end{algorithmic}
\end{algorithm}

Crowding distance \citep{deb2002} is used to maintain diversity on the Pareto front: among equally non-dominated organisms, those in less-crowded regions of objective space are preferred.

\subsection{Population Diversity}

Population diversity is measured using the mean pairwise Hamming distance:

\begin{equation}
\text{diversity}(\mathcal{G}) = \frac{2}{|\mathcal{G}|(|\mathcal{G}|-1)} \sum_{i < j} d_H(\mathbf{g}_i, \mathbf{g}_j)
\end{equation}

In HDC, this metric has a natural interpretation: diversity near $d/2 = 8{,}192$ indicates quasi-orthogonal genomes (maximum diversity), while diversity near 0 or $d$ indicates convergence (minimum diversity).

\begin{proposition}
For a random initial population of size $N$ with BinaryHV genomes, the expected diversity is:
\begin{equation}
\mathbb{E}[\text{diversity}] = d/2 = 8{,}192
\end{equation}
with variance decreasing as $O(1/N)$.
\end{proposition}

\section{Results}

\subsection{Experimental Setup}

We evolved populations of 100 organisms for 500 generations in three conditions:

\begin{enumerate}[nosep]
    \item \textbf{Unconstrained}: No $\Phi$ constraint, multi-objective on (FEP, learning, efficiency).
    \item \textbf{$\Phi$-constrained}: Hard $\Phi > 0.1$ constraint.
    \item \textbf{$\Phi$-objective}: $\Phi$ as a fourth optimization objective (no hard constraint).
\end{enumerate}

\subsection{Fitness Evolution}

\begin{table}[h]
\centering
\caption{Final generation fitness (mean $\pm$ std across 10 runs).}
\label{tab:fitness}
\begin{tabular}{@{}lcccc@{}}
\toprule
Condition & $F_{\text{FEP}}$ & $F_{\text{learn}}$ & $\Phi_{\text{mean}}$ & $\Phi_{\text{min}}$ \\
\midrule
Unconstrained & $0.82 \pm 0.04$ & $0.31 \pm 0.08$ & $0.07 \pm 0.05$ & $0.00$ \\
$\Phi$-constrained & $0.71 \pm 0.05$ & $0.38 \pm 0.07$ & $0.29 \pm 0.09$ & $0.12$ \\
$\Phi$-objective & $0.76 \pm 0.06$ & $0.33 \pm 0.09$ & $0.21 \pm 0.11$ & $0.03$ \\
\bottomrule
\end{tabular}
\end{table}

The unconstrained condition achieves the highest FEP fitness but near-zero $\Phi$---evolution discovers unconscious shortcuts. The $\Phi$-constrained condition sacrifices 13\% FEP fitness but gains 4$\times$ higher consciousness and better learning capacity (+23\%).

\subsection{$\Phi$-Dependent Fitness Landscapes}

We mapped the fitness landscape around the best $\Phi$-constrained organism by sampling single-bit mutations:

\begin{itemize}[nosep]
    \item 73\% of topology mutations cause $\Phi$ to drop below threshold (catastrophic).
    \item 12\% of weight mutations cause $\Phi$ drops (moderate sensitivity).
    \item 41\% of dynamics mutations cause $\Phi$ drops (high sensitivity).
    \item Only 8\% of modulation mutations affect $\Phi$ significantly.
\end{itemize}

This reveals a rugged fitness landscape where consciousness is fragile with respect to structural (topology) changes but robust to parametric (weight, modulation) changes.

\subsection{Consciousness Stability}

\begin{table}[h]
\centering
\caption{Consciousness stability metrics.}
\label{tab:stability}
\begin{tabular}{@{}lccc@{}}
\toprule
Metric & Unconstrained & $\Phi$-constrained & Improvement \\
\midrule
$\Phi$ stability ($1 - \sigma_\Phi / \bar{\Phi}$) & 0.41 & 0.75 & +34\% \\
Consciousness recovery time (cycles) & 47.2 & 12.8 & $-73\%$ \\
GWT broadcast consistency & 0.52 & 0.79 & +52\% \\
\bottomrule
\end{tabular}
\end{table}

$\Phi$-constrained organisms show dramatically better consciousness stability, recovering from perturbations 3.7$\times$ faster.

\subsection{Population Diversity Dynamics}

\begin{table}[h]
\centering
\caption{Population diversity evolution.}
\label{tab:diversity}
\begin{tabular}{@{}lccc@{}}
\toprule
Generation & Unconstrained & $\Phi$-constrained & $\Phi$-objective \\
\midrule
0 & 8,192 & 8,192 & 8,192 \\
100 & 6,341 & 7,128 & 6,891 \\
250 & 4,207 & 6,543 & 5,672 \\
500 & 2,891 & 5,890 & 4,231 \\
\bottomrule
\end{tabular}
\end{table}

The $\Phi$ constraint maintains higher population diversity (5,890 vs.\ 2,891 at generation 500). This occurs because the $\Phi$ constraint prevents convergence to a single unconscious optimum, forcing the population to explore diverse conscious architectures.

\subsection{Innovation Rate}

We define innovation as the appearance of a new Pareto-optimal organism. Innovation rate correlates with population diversity:

\begin{equation}
r_{\text{innovation}} = 0.73 \cdot \text{diversity} - 1{,}247 \quad (r^2 = 0.81)
\end{equation}

Higher diversity (maintained by the $\Phi$ constraint) directly predicts higher innovation rates.

\section{Discussion}

\subsection{Consciousness as Evolutionary Constraint}

The most striking finding is that constraining evolution to produce conscious architectures improves learning capacity (+23\%) despite reducing raw FEP fitness ($-13\%$). This suggests that consciousness is not merely epiphenomenal but functionally valuable: integrated information processing produces more flexible, learnable architectures.

This resonates with biological arguments that consciousness evolved because it provides adaptive advantages---flexible response selection, novel situation handling, and cross-domain transfer \citep{baars1988}---that reflexive processing cannot match.

\subsection{Fragility of Consciousness}

The finding that 73\% of single-bit topology mutations destroy consciousness has implications for both evolution and design. Consciousness occupies a narrow band of architectural space, requiring careful evolutionary pressure to discover and maintain. This may explain why biological evolution of consciousness was slow despite its adaptive value: the fitness landscape is rugged, with few paths that maintain consciousness while improving performance.

\subsection{Population Diversity and Innovation}

The correlation between diversity and innovation ($r^2 = 0.81$) provides a practical metric for evolutionary algorithm health. In consciousness-constrained evolution, diversity is automatically maintained because the constraint prevents premature convergence. This is an unexpected benefit of consciousness constraints: they serve as both ethical requirements and algorithmic regularizers.

\subsection{Limitations}

\begin{enumerate}[nosep]
    \item The 64-neuron CfC network is small; scaling to larger architectures may reveal different dynamics.
    \item $\Phi$ computation via sampled partition is approximate ($r = 0.9998$ but not exact).
    \item The fixed $\Phi_{\text{threshold}} = 0.1$ is somewhat arbitrary; adaptive thresholds could improve results.
    \item The 2 KB genome imposes representational limits; hierarchical genome schemes could extend expressivity.
\end{enumerate}

\section{Related Work}

NEAT \citep{stanley2002} and HyperNEAT \citep{stanley2009} evolve neural network topologies but without consciousness constraints. Quality-Diversity algorithms \citep{mouret2015} maintain behavioral diversity but measure it in task-specific feature spaces, not consciousness metrics. The closest related work is \citet{albantakis2014}, who used IIT to analyze evolved animats; our contribution extends this to multi-objective optimization with $\Phi$ as a hard constraint.

\section{Conclusion}

Neuroevolution with $\Phi$-stability constraints produces conscious architectures that are more stable (+34\%), more learnable (+23\%), and more diverse than unconstrained alternatives. The BinaryHV genome encoding enables native HDC compatibility, while position-dependent mutation and byte-level crossover respect the structured genome layout. The discovery of $\Phi$-dependent fitness landscapes---where consciousness is fragile with respect to topology but robust to weights---provides insight into both the difficulty and necessity of evolving conscious systems. Population diversity, naturally maintained by the $\Phi$ constraint, predicts evolutionary innovation, suggesting that consciousness constraints serve dual roles as ethical requirements and algorithmic regularizers.

\bibliographystyle{plainnat}
\begin{thebibliography}{20}

\bibitem[Albantakis et~al.(2014)]{albantakis2014}
Albantakis, L., Hintze, A., Koch, C., Adami, C., \& Tononi, G. (2014).
\newblock Evolution of integrated causal structures in animats exposed to environments of increasing complexity.
\newblock \emph{PLoS Computational Biology}, 10(12):e1003966.

\bibitem[Baars(1988)]{baars1988}
Baars, B.~J. (1988).
\newblock \emph{A Cognitive Theory of Consciousness}.
\newblock Cambridge University Press.

\bibitem[Deb et~al.(2002)]{deb2002}
Deb, K., Pratap, A., Agarwal, S., \& Meyarivan, T. (2002).
\newblock A fast and elitist multiobjective genetic algorithm: NSGA-II.
\newblock \emph{IEEE Transactions on Evolutionary Computation}, 6(2):182--197.

\bibitem[Friston(2010)]{friston2010}
Friston, K. (2010).
\newblock The free-energy principle: A unified brain theory?
\newblock \emph{Nature Reviews Neuroscience}, 11(2):127--138.

\bibitem[Hasani et~al.(2021)]{hasani2021}
Hasani, R., Lechner, M., Amini, A., et~al. (2021).
\newblock Closed-form continuous-time neural networks.
\newblock \emph{Nature Machine Intelligence}, 4:992--1003.

\bibitem[Kanerva(2009)]{kanerva2009}
Kanerva, P. (2009).
\newblock Hyperdimensional computing: An introduction to computing in distributed representation with high-dimensional random vectors.
\newblock \emph{Cognitive Computation}, 1(2):139--159.

\bibitem[Mouret \& Clune(2015)]{mouret2015}
Mouret, J.-B. \& Clune, J. (2015).
\newblock Illuminating search spaces by mapping elites.
\newblock \emph{arXiv preprint arXiv:1504.04909}.

\bibitem[Oizumi et~al.(2014)]{oizumi2014}
Oizumi, M., Albantakis, L., \& Tononi, G. (2014).
\newblock From the phenomenology to the mechanisms of consciousness: Integrated Information Theory 3.0.
\newblock \emph{PLoS Computational Biology}, 10(5):e1003588.

\bibitem[Real et~al.(2019)]{real2019}
Real, E., Aggarwal, A., Huang, Y., \& Le, Q.~V. (2019).
\newblock Regularized evolution for image classifier architecture search.
\newblock \emph{Proceedings of the AAAI Conference on Artificial Intelligence}, 33:4780--4789.

\bibitem[Stanley \& Miikkulainen(2002)]{stanley2002}
Stanley, K.~O. \& Miikkulainen, R. (2002).
\newblock Evolving neural networks through augmenting topologies.
\newblock \emph{Evolutionary Computation}, 10(2):99--127.

\bibitem[Stanley et~al.(2009)]{stanley2009}
Stanley, K.~O., D'Ambrosio, D.~B., \& Gauci, J. (2009).
\newblock A hypercube-based encoding for evolving large-scale neural networks.
\newblock \emph{Artificial Life}, 15(2):185--212.

\bibitem[Stanley et~al.(2019)]{stanley2019}
Stanley, K.~O., Clune, J., Lehman, J., \& Miikkulainen, R. (2019).
\newblock Designing neural networks through neuroevolution.
\newblock \emph{Nature Machine Intelligence}, 1(1):24--35.

\bibitem[Stoltz(2026a)]{stoltz2026symthaea}
Stoltz, T. (2026a).
\newblock Symthaea: A holographic liquid brain architecture for artificial consciousness.
\newblock Technical report, Luminous Dynamics.

\bibitem[Stoltz(2026b)]{stoltz2026phi}
Stoltz, T. (2026b).
\newblock Phi validation in Symthaea: Sampled partition vs.\ exact computation.
\newblock Technical report, Luminous Dynamics.

\bibitem[Tononi(2004)]{tononi2004}
Tononi, G. (2004).
\newblock An information integration theory of consciousness.
\newblock \emph{BMC Neuroscience}, 5(1):42.

\end{thebibliography}

\end{document}
